% ID: SEL-D0026-OB-E03
% ESTADO: cotejada. Segunda pasada visual completa y borrador compilado revisados.
% Fuente: Selectividad/matematicas_aplicadas/suplemente_septiembre_EXAMEN 4-ANDALUCÍA.pdf
% SHA-256: c41ea8ffe46f2de3a4d7f82e43def97322e29653182ad3df54a57e2ead2beb19
% Página: 2; opción B; ejercicio original 3.

De los alumnos que se presentaron a las pruebas de selectividad de una provincia, 1150 se examinaron de Geografía; de estos, 598 eligieron la opción A. Se sabe que aprobaron esa asignatura el 78\% de los que eligieron la opción A y el 74\% de los que eligieron la opción B. Se ha escogido al azar uno de los alumnos que se examinaron de Geografía.
\begin{enumerate}[label=\alph*)]
\item \textbf{(1.5 puntos)} \textquestiondown{}Cuál es la probabilidad de que este alumno haya aprobado esta asignatura?
\item \textbf{(1 punto)} Si se sabe que este alumno ha aprobado Geografía, \textquestiondown{}cuál es la probabilidad de que haya elegido la opción A?
\end{enumerate}
